% Chapter 1: Introduction

\chapter{Introduction}\label{ch:introduction}

\section{Background and Motivation}\label{sec:background}

The reliability of industrial robots in assembly, grasping, and human--robot collaboration tasks depends critically on the accuracy of their joint encoders. Whether at the level of low-level joint servoing, Cartesian impedance control, or high-level task planning, both the control law and the observation module take the encoder reading $\vect{q}_{\text{measured}}$ as feedback input, and map joint space to task space through forward kinematics and the Jacobian~\cite{siciliano2010robotics}. In real-world deployment, however, a systematic bias $\vect{b}$ between the encoder reading and the true joint angle is pervasive, such that $\vect{q}_{\text{measured}} = \vect{q}_{\text{true}} + \vect{b}$.

This bias can be introduced by a variety of mechanisms---factory calibration residuals, temperature drift, replacement of components without recalibration, misalignment after a collision, electromagnetic interference, among others (a detailed taxonomy and typical magnitudes are given in \cref{sec:bias_sources}). Among these, the ``replacement without recalibration'' and ``post-collision misalignment'' categories can reach magnitudes above $0.1$\,rad, which is sufficient to significantly degrade an otherwise reliable policy when it executes fine manipulation tasks. To cope with such bias, industrial practice typically adopts two approaches: \emph{offline recalibration}~\cite{kana2022fast,gong2000nongeometric} (high in precision, but requiring an interruption of production that disrupts the line takt time) and \emph{traditional fault-tolerant control} (FDI/FTC)~\cite{blanke2016diagnosis}, which relies on an explicit fault model and aims at ``maintaining stability'' rather than ``completing the task,'' and is therefore often powerless against a problem such as a latent bias of unknown magnitude. This thesis takes a path different from both: \textbf{without interrupting production, the policy is exposed to a variety of biases already during training, so that it can adapt online to an unknown calibration bias.} A comparison and positioning of several natural alternatives to this path---such as explicitly identifying the bias and then compensating for it (which still requires an independent external observation and an offline recomputation of the Jacobian), or meta-RL (which, under the assumption of a constant latent variable, degenerates into single-task domain randomization)---is given in \cref{ch:related_work}. This idea calls for a re-examination of the formal structure of the control problem (the bias as a latent variable degrades the MDP to a POMDP), the observation design, the training paradigm, and the guarantee of human--robot safety; these questions constitute the research theme of this thesis.

\section{Problem Statement}\label{sec:problem_statement}

Once an unknown bias $\vect{b}$ is introduced, the original MDP breaks down at three interrelated levels.

\paragraph{Partial observability} A single bias source produces effects along three independent paths---an initial position offset (the steady-state result of the PD control law deviates from $\vect{q}_{\text{home}}$), a control bias (the task-space torque is computed from an incorrect Jacobian), and a perception bias (both the joint reading and the forward-kinematics end-effector pose are corrupted simultaneously). All three effects make it impossible for the agent to directly observe the true joint state, and the control problem degrades from an MDP to a POMDP (\cref{ch:problem_formulation} further proves that the bias is non-identifiable under the basic observation, and becomes approximately identifiable once an independent external position observation is introduced).

\paragraph{Deployment constraints} The visual sim-to-real gap causes pure-simulation training followed by direct transfer to fail on the real robot; real-robot data collection is costly, with the number of available demonstrations on the order of $50$--$200$. \textbf{Under the real-robot conditions of this thesis} ($10$\,Hz control frequency, a single episode $\le 30$\,s, on the order of $50$ demonstrations), under a sparse reward the SAC online stage exhibits a large drift once the actor is unfrozen (a detailed post-mortem is given in \cref{subsec:hilserl_failure}). Together, these constraints determine a dual-end paradigm of ``simulation-side validation plus real-robot imitation deployment'': the simulation side uses RL to validate theoretical feasibility, and the real-robot side uses imitation learning (IL; including behavior cloning, BC, and human-in-the-loop intervention correction) to carry out deployment.

\paragraph{Human--robot collaborative safety} Real-robot deployment involves a human hand intermittently entering the workspace (operator intervention during HIL training, and a worker approaching during production deployment), which requires a backup avoidance policy that cooperates with the task policy, taking over control when a distance threshold is triggered and handing it back once the avoidance is complete. The backup policy itself must be trained without introducing additional unsafe factors; this thesis adopts a scheme of pure-simulation SAC plus domain randomization (the safety constraint during HIL training is likewise provided by this backup policy, interrelated with the previous class of constraints).

\section{Contributions}\label{sec:contributions}

The core contributions of this thesis revolve around its title, ``Hierarchical Imitation and Reinforcement Learning for Fault-Tolerant Manipulation,'' and fall into three points.

\textbf{Contribution 1: POMDP formalization of encoder bias and observability analysis} (\cref{ch:problem_formulation}). This thesis presents a one-source-three-effect causal-chain model and a seven-tuple POMDP formalization, and proves through rigorous propositions that $(\vect{q}_{\text{true}}, \vect{b})$ is non-identifiable under the basic observation (\cref{prop:non_identifiable}) and becomes approximately identifiable once an independent external position observation is introduced (\cref{prop:identifiable}). This is \emph{identifiability in the information-theoretic sense}, and does not directly imply the learnability of any specific policy class; \cref{subsec:identifiability_vs_solvability} further clarifies that ``identifiability $\neq$ solvability''---even when the information is sufficient, the optimal policy must still output a compensating action conditioned on $\vect{b}$, which in turn establishes the necessity of ``randomized-bias training.''

\textbf{Contribution 2: A fault-tolerant manipulation framework of hierarchical imitation and reinforcement learning} (\cref{ch:task_policy}--\cref{ch:real_chapter}). This thesis designs a task $\oplus$ backup dual-policy hierarchical architecture together with a dual-end paradigm split into sim/real, comprising:

\begin{enumerate}
  \item \textbf{A shared method layer} (\cref{ch:task_policy})---observation design, randomized-bias training, a shared network architecture, and a sparse reward;
  \item \textbf{Simulation-side task-policy training} (\cref{ch:sim_chapter})---the HIL-SERL framework validates the observability propositions (the H1--H3 experimental loop);
  \item \textbf{Real-robot-side task-policy deployment} (\cref{ch:real_chapter})---BC plus HG-DAgger imitation iteration, which avoids the actor-drift problem of SAC in the small-data real-robot setting;
  \item \textbf{Backup-policy simulation training and real-robot deployment}---pure-simulation SAC plus DR training (\cref{sec:backup_sim_train}), with a hard switch between the task policy and the backup policy carried out on the real robot by the three-state FSM of the Hierarchical Supervisor (\cref{subsec:supervisor}).
\end{enumerate}

\noindent The paradigm choice is formalized by a three-axis decision principle of (observation modality $\times$ data scale $\times$ safety constraint) (see \cref{sec:task_method_motivation} and \cref{sec:backup_paradigm} for details); the task and backup policies are two instances of this decision principle on the system of this thesis.

\textbf{Contribution 3: End-to-end validation of three real-robot tasks under injected encoder bias} (\cref{sec:real_experiments}). This thesis implements a B$+$D dual-injection-point encoder-bias injection (\cref{subsec:bd_injection}), and designs a four-column joint success-rate matrix evaluation (a $2\times2$ ``training bias $\times$ test bias'' factorial: train w/o bias \& test w/o bias, train w/o bias \& test w/ bias, train w/ bias \& test w/o bias, train w/ bias \& test w/ bias) on the three real-robot tasks \texttt{pickup}, \texttt{wipe}, and \texttt{pickandplace}, which constitutes the core empirical thread of this thesis (the experimental design and partial results are given in \cref{sec:real_experiments}, and the full matrix continues to be filled in as real-robot data accumulate). It also presents a methodology-level negative result of the HIL-SERL real-robot failure (see \cref{subsec:hilserl_failure}), which provides empirical support for the paradigm choice of switching to BC plus HG-DAgger on the real robot.

\section{Thesis Organization}\label{sec:organization}

This thesis consists of seven chapters, organized in the following order:

\begin{itemize}
  \item \textbf{Chapter 1, Introduction}: research background, problem statement, contribution claims, and organizational structure.
  \item \cref{ch:related_work} (Related Work): reviews the literature on FDI/FTC, POMDP solving, imitation learning and human-in-the-loop methods, and safe RL, and positions the methodological choices of this thesis.
  \item \cref{ch:problem_formulation} (Encoder Bias Modeling and POMDP Analysis): the theoretical core---the one-source-three-effect causal chain, the POMDP formalization, and the observability propositions.
  \item \cref{ch:task_policy} (Task Policy: Method Design and Shared Network Architecture): the method layer shared by simulation and the real robot.
  \item \cref{ch:sim_chapter} (Simulation: Policy Training and Observability Validation): the two main threads on the simulation side (the HIL-SERL task policy and the SAC$+$DR backup policy) and the H1--H3 experimental loop.
  \item \cref{ch:real_chapter} (Real Robot: Imitation Deployment and End-to-End Validation): the full real-robot deployment stack and the three-task joint success-rate matrix.
  \item \cref{ch:conclusion} (Conclusion and Outlook): summarizes the contributions and gives future work (including the two-step ablation of real-robot bias robustness).
\end{itemize}
